\chapter{Probability Spaces and \texorpdfstring{$\sigma$-}{sigma-}algebras}

\section{Measurable Spaces}

The foundation of rigorous probability theory lies in measure theory. Before we can talk about probabilities of events, we must strictly define what an ``event'' is. Not every subset of a sample space can consistently be assigned a probability, especially when dealing with uncountably infinite sample spaces like the real line $\mathbb{R}$.

\begin{definition}[\texorpdfstring{$\sigma$-}{sigma-}algebra]
Let $\Omega$ be a non-empty set (the sample space). A collection $\mathcal{F}$ of subsets of $\Omega$ (i.e., $\mathcal{F} \subseteq 2^{\Omega}$) is called a \textbf{\texorpdfstring{$\sigma$-}{sigma-}algebra} (or \texorpdfstring{$\sigma$-}{sigma-}field) on $\Omega$ if it satisfies the following three axioms:
\begin{enumerate}
    \item \textbf{Empty set}: $\emptyset \in \mathcal{F}$.
    \item \textbf{Complement}: If $A \in \mathcal{F}$, then its complement $A^c := \Omega \setminus A \in \mathcal{F}$.
    \item \textbf{Countable Union}: If $(A_n)_{n \geq 1} \subset \mathcal{F}$ is a sequence of sets in $\mathcal{F}$, then their union is also in $\mathcal{F}$:
    \[
    \bigcup_{n \geq 1} A_n \in \mathcal{F}.
    \]
\end{enumerate}
The pair $(\Omega, \mathcal{F})$ is called a \textbf{measurable space}. The elements of $\mathcal{F}$ are called \textbf{measurable sets} or, in the context of probability, \textbf{events}.
\end{definition}

\begin{remark}
From these axioms, we can deduce other properties:
\begin{itemize}
    \item The whole space $\Omega$ must be in $\mathcal{F}$ because $\Omega = \emptyset^c$.
    \item \texorpdfstring{$\sigma$-}{sigma-}algebras are closed under countable intersections. By De Morgan's laws:
    \[
    \bigcap_{n \geq 1} A_n = \left( \bigcup_{n \geq 1} A_n^c \right)^c \in \mathcal{F}.
    \]
\end{itemize}
\end{remark}

\subsection{Examples of \texorpdfstring{$\sigma$-}{sigma-}algebras}
\begin{enumerate}
    \item \textbf{Trivial \texorpdfstring{$\sigma$-}{sigma-}algebra}: $\mathcal{F} = \{ \emptyset, \Omega \}$. This is the smallest possible \texorpdfstring{$\sigma$-}{sigma-}algebra. It corresponds to having no information about the outcome of an experiment other than it happened.
    \item \textbf{Power set}: $\mathcal{F} = 2^{\Omega} = \{ A : A \subset \Omega \}$. This is the largest possible \texorpdfstring{$\sigma$-}{sigma-}algebra, containing every possible subset. While useful for discrete spaces (like rolling a die), it is often too large for continuous spaces like $\mathbb{R}$ to define a consistent measure (Vitali sets).
    \item \textbf{Partition-generated}: Let $(A_i)_{i \in I}$ be a countable partition of $\Omega$, meaning $\Omega = \bigsqcup_{i \in I} A_i$ (disjoint union). The \texorpdfstring{$\sigma$-}{sigma-}algebra generated by this partition consists of all possible unions of the sets $A_i$:
    \[
    \mathcal{F} = \left\{ \bigcup_{j \in J} A_j : J \subseteq I \right\}.
    \]
    This structure is very important in finance when modeling information flow in discrete time. If we know which partition element $\omega$ falls into, our information is described by this \texorpdfstring{$\sigma$-}{sigma-}algebra.
\end{enumerate}

\section{Generated \texorpdfstring{$\sigma$-}{sigma-}algebras}

Often provided with a collection of sets of interest, we want to construct the smallest \texorpdfstring{$\sigma$-}{sigma-}algebra that allows us to measure them.

\begin{definition}[Generated \texorpdfstring{$\sigma$-}{sigma-}algebra]
Let $\mathcal{C} \subseteq 2^{\Omega}$ be any collection of subsets of $\Omega$. The \texorpdfstring{$\sigma$-}{sigma-}algebra generated by $\mathcal{C}$, denoted $\sigma(\mathcal{C})$, is the intersection of all \texorpdfstring{$\sigma$-}{sigma-}algebras containing $\mathcal{C}$:
\[
\sigma(\mathcal{C}) := \bigcap_{\mathcal{F} \supset \mathcal{C}, \mathcal{F} \text{ is a } \sigma\text{-algebra}} \mathcal{F}.
\]
It is the \textbf{smallest} \texorpdfstring{$\sigma$-}{sigma-}algebra containing $\mathcal{C}$.
\end{definition}

\subsection{Borel \texorpdfstring{$\sigma$-}{sigma-}algebra}
In topological spaces, we want our algebraic structure (events) to be compatible with the topological structure (open sets, convergence).

\begin{definition}[Borel \texorpdfstring{$\sigma$-}{sigma-}algebra]
Let $(\Omega, \tau)$ be a topological space. The \textbf{Borel \texorpdfstring{$\sigma$-}{sigma-}algebra}, denoted $\mathcal{B}(\Omega)$, is the \texorpdfstring{$\sigma$-}{sigma-}algebra generated by the family of open sets $\mathcal{O}$ in $\Omega$.
\end{definition}

\subsubsection{Case of $\mathbb{R}$ and $\mathbb{R}^d$}
\begin{itemize}
    \item On $\mathbb{R}$, $\mathcal{B}(\mathbb{R})$ is generated by open intervals $]a, b[$. It is the smallest family of sets containing all intervals that is closed under countable operations.
    \item On $\mathbb{R}^d$, $\mathcal{B}(\mathbb{R}^d)$ is generated by open rectangles (or ``hyper-rectangles''):
    \[
    \prod_{i=1}^d ]a_i, b_i[ = \{ \omega = (\omega_1, \dots, \omega_d) \in \mathbb{R}^d : a_i < \omega_i < b_i, \forall i \}.
    \]
\end{itemize}
Note that while $\mathcal{B}(\mathbb{R}^d)$ contains almost any set we can construct explicitly, there do exist non-Borel sets (though they are pathological).

\subsection{Product \texorpdfstring{$\sigma$-}{sigma-}algebra}
When dealing with multiple random variables or stochastic processes, we work with Cartesian products of spaces.

Let $((\Omega_n, \mathcal{F}_n))_{n \geq 1}$ be a sequence of measurable spaces.
\begin{itemize}
    \item \textbf{Finite Product}: For a finite number of spaces $\Omega_1 \times \dots \times \Omega_n$, the product \texorpdfstring{$\sigma$-}{sigma-}algebra $\mathcal{F}_1 \otimes \dots \otimes \mathcal{F}_n$ is generated by ``rectangles'' of the form:
    \[
    A_1 \times \dots \times A_n, \quad \text{where } A_i \in \mathcal{F}_i.
    \]
    Important property: $\mathcal{B}(\mathbb{R}) \otimes \dots \otimes \mathcal{B}(\mathbb{R}) = \mathcal{B}(\mathbb{R}^n)$.
    
    \item \textbf{Infinite Product}: For the infinite product space $\prod_{n \geq 1} \Omega_n$, the product $\sigma$-algebra $\bigotimes_{n \geq 1} \mathcal{F}_n$ is generated by \textbf{cylinders} with a finite basis. A measurable cylinder is a set of the form:
    \[
    C = \{ \omega \in \prod \Omega_n : \omega_{i_1} \in A_{i_1}, \dots, \omega_{i_k} \in A_{i_k} \}
    \]
    where only a finite number of coordinates $i_1, \dots, i_k$ are restricted to subsets $A_{i_j} \in \mathcal{F}_{i_j}$, and all other coordinates are free (i.e., in $\Omega_m$).
    This definition is crucial for defining stochastic processes as random variables taking values in spaces of functions sequences.
\end{itemize}

\section{Probability Measures}

\begin{definition}[Measure]
A map $\mu: \mathcal{F} \to [0, \infty]$ is called a \textbf{measure} if:
\begin{enumerate}
    \item $\mu(\emptyset) = 0$.
    \item \textbf{$\sigma$-additivity}: For any sequence $(A_n)_{n \geq 1} \subset \mathcal{F}$ of \textbf{pairwise disjoint} sets,
    \[
    \mu\left( \bigcup_{n \geq 1} A_n \right) = \sum_{n \geq 1} \mu(A_n).
    \]
\end{enumerate}
\end{definition}

\begin{definition}[Probability Measure]
A measure $\mathbb{P}$ is a \textbf{probability measure} if it satisfies the normalization axiom:
\[
\mathbb{P}(\Omega) = 1.
\]
The triplet $(\Omega, \mathcal{F}, \mathbb{P})$ is called a \textbf{probability space}.
\end{definition}

We denote the set of all probability measures on $(\Omega, \mathcal{F})$ as $\mathcal{P}(\Omega, \mathcal{F})$.

\section{Support of a Measure}

The concept of support formalizes the idea of ``where the probability mass lives''.

\begin{definition}[Support]
The \textbf{support} of a measure $\mu$, denoted $\text{supp}(\mu)$, is the closed set in $\Omega$ such that:
\begin{enumerate}
    \item The measure of its complement is zero: $\mu(\Omega \setminus \text{supp}(\mu)) = 0$.
    \item For any open set $G$ such that $G \cap \text{supp}(\mu) \neq \emptyset$, we have $\mu(G \cap \text{supp}(\mu)) > 0$.
\end{enumerate}
Essentially, it is the smallest closed set that carries all the mass of $\mu$. If measurement values fall outside the support, it is an ``impossible'' event (probability 0).
\end{definition}
For a Polish space (separable, completely metrizable, like $\mathbb{R}^d$), the support is unique and given by:
\[
\text{supp}(\mu) = \Omega \setminus \bigcup \{ V : V \text{ is open and } \mu(V) = 0 \}.
\]
